\chapter{Ureteroscopy}

\section*{Overview}
Ureteroscopy involves passing a thin, rigid or flexible fiberoptic scope through the urethra and bladder into the ureter to diagnose and treat ureteral stones or tumors. The exact approach, setting, and preparation depend on the indication, anatomy, and the patient's health.

\section*{Alternative Names}
URS, Ureteral stone scope

\section*{Principle}
A thin endoscope is passed through the urethra, bladder, and ureter, giving direct visualization of the ureter and renal pelvis. Lasers and baskets passed through the scope fragment and extract stones, and strictures can be dilated or incised.
\section*{History}
Hugh Hampton Young performed the first ureteroscopy in a patient with valves in 1912. The technique advanced in the 1980s with flexible scopes and laser fibers.

\section*{Why It Is Done}
Ordered to treat ureteral stones (especially in the middle or lower ureter) that cannot be treated by lithotripsy, or to diagnose ureteral strictures or tumors.

\section*{Is This a Primary or Secondary Test or Procedure}
Primary endoscopic treatment for ureteral stones and upper tract urothelial carcinoma biopsy.

\section*{How It Is Performed}
Performed under general anesthesia. The ureteroscope is guided through the urinary tract. A holmium laser fiber is used to fragment the stone, and fragments are removed with a basket. A ureteral stent is often placed.

\section*{Preparation}
Fasting depends on the anesthesia plan. The team evaluates for urinary infection and reviews anticoagulants, antiplatelet medicines, allergies, and kidney function. Do not stop blood-thinning medicines unless the prescribing and procedural teams provide a coordinated plan.

\section*{Contraindications and Precautions}
Untreated urinary tract infection, or severe urethral stricture.

\section*{Risks}
Ureteral injury (perforation, avulsion), hematuria, UTI, or stent pain/cramping.

\section*{Aftercare and Recovery}
Drink plenty of water. If a ureteral stent was placed, you will feel a frequent urge to urinate and mild flank pain during urination. The stent is removed in 3-10 days.

\section*{Outcome and Follow-up}
Direct visualization confirms complete stone clearance and normal ureteral mucosa.

\section*{Expected Results}
Complete clearance of ureteral stone; normal ureteral mucosa; patent ureteral lumen.

\section*{Follow-up}
Stent removal; follow-up ultrasound or X-ray in 2-4 weeks to check for hydronephrosis.

\section*{When to Seek Medical Care}
Follow the procedure team's specific instructions. Seek urgent medical assessment for severe or worsening pain, uncontrolled bleeding, fever or signs of infection, trouble breathing, chest pain, fainting, new weakness or confusion, inability to urinate, or any rapidly worsening symptom. The appropriate warning signs depend on the procedure and the patient's medical condition.

\section*{References}
\begin{enumerate}
  \item MedlinePlus. Ureteroscopy. U.S. National Library of Medicine; 2025.
  \item Current Medical Diagnosis and Treatment. McGraw-Hill; 2025.
\end{enumerate}

\clearpage
